% Figure: the spiral heat exchanger, two metal sheets wound into a pair of concentric
% spiral channels so that the hot stream in one channel runs counter-current to the
% cold stream in the neighbouring channel across the whole wound length; the single
% long curved passage per side is what gives the unit its self-scouring, plug-tolerant
% character.
\begin{figure}[htbp]
\centering
\begin{tikzpicture}[scale=1.0]
% draw two interleaved Archimedean spirals as sampled polar curves
% hot channel spiral (red), cold channel spiral (blue), offset by half a turn
\draw[engred, very thick, domain=0:1080, samples=200, smooth]
  plot ({(0.35+0.020*\x)*cos(\x)}, {(0.35+0.020*\x)*sin(\x)});
\draw[engblue, very thick, domain=0:1080, samples=200, smooth]
  plot ({(0.55+0.020*\x)*cos(\x+180)}, {(0.55+0.020*\x)*sin(\x+180)});
% outer shell circle
\draw[enggray, thick] (0,0) circle (23mm);
% hot inlet arrow at the outer edge (feeds the outermost part of the red spiral)
\draw[engred, -{Latex[length=2.2mm]}, very thick] (3.2,0.6) -- (2.15,0.35);
\node[engred, font=\scriptsize, anchor=west] at (3.2,0.7) {hot in};
% hot outlet leaves at the centre
\draw[engred, -{Latex[length=2.2mm]}, very thick] (0.15,0.0) -- (0.7,-0.9);
\node[engred, font=\scriptsize, anchor=north] at (0.75,-0.95) {hot out (core)};
% cold inlet at the centre
\draw[engblue, -{Latex[length=2.2mm]}, very thick] (-0.9,0.9) -- (-0.35,0.35);
\node[engblue, font=\scriptsize, anchor=east] at (-0.95,0.95) {cold in (core)};
% cold outlet at the outer edge
\draw[engblue, -{Latex[length=2.2mm]}, very thick] (-2.1,-0.4) -- (-3.1,-0.7);
\node[engblue, font=\scriptsize, anchor=east] at (-3.15,-0.75) {cold out};
% label the counter-current pair
\node[enggray, font=\scriptsize, anchor=north] at (0,-2.6) {two sheets wound into interleaved spiral channels; one curved pass per side};
\end{tikzpicture}
\caption[Spiral heat exchanger channel arrangement]{The spiral heat exchanger, formed
by winding two metal sheets into a pair of interleaved spiral channels. The hot
stream enters at the outer edge and spirals inward to leave at the core, while the
cold stream enters at the core and spirals outward, so the two run counter-current
across the whole wound length with a single long curved passage on each side.
The single passage per side is the unit's defining feature: there is no bundle of
parallel tubes for a deposit to block selectively, and the continuously curving flow
raises the local velocity and the turbulence at the channel walls, so a spiral
exchanger scours its own surfaces and tolerates slurries and fouling fluids that
would plug a shell-and-tube unit. The price is a fixed geometry that cannot be
re-bundled and a pressure envelope set by the wound sheets rather than by a code
shell, which is why the spiral is chosen for dirty, viscous, and fibrous services
rather than for high-pressure clean duty.}
\label{fig:vol04ch06:spiral-exchanger}
\end{figure}
